% page 156 of the facsimile (image p-155 of the scan)
% block 160.0 x 229.0 mm ; left margin 8.0 mm
\pgc{-12.700mm}{0.000mm}{
\l{\cel{5}146.}
\l{}
\l{}
\l{\sou{estetiko.}\cel{2}Parto di filozofio, qua traktas lo bela en naturo}
\l{\cel{3}ed en arto. - DEFIRS.}
\l{}
\l{\sou{estimar}. (\sou{trans.}) Opinionar favoroze pri la valoro (etikala,}
\l{\cel{3}morala, karakterala) di persono. - DEFIS.}
\l{}
\l{\sou{estivar}. (\sou{netrans}.) (\sou{pri la bestio-trupi)}. Vivar dum somero}
\l{\cel{3}en la pastureyi di la monti. - DEFIRS.}
\l{}
\l{\sou{estompar}.\cel{2}(\sou{trans.}) I. Dilutar la traci de krayono e pastelo,}
\l{\cel{3}sur desegnuro, per rolero felo o papera, di qua la extre-}
\l{\cel{3}maji esas ordinare pinta. - II. Kovrar ye nuanco subtila.}
\l{\cel{3}- DEF.}
\l{}
\l{\sou{-estr-}.\cel{2}Sufixo qua signifikas \textquotedbl{}chefo di\textquotedbl{}, \textquotedbl{}mastro di\textquotedbl{}.}
\l{}
\l{\sou{estrado}. Parto elevita super la planko-sulo di chambro, di}
\l{\cel{3}chambrego, di salono, sur qua onu lokizas lito, trono,}
\l{\cel{3}katedro, e c. - DEFRS.}
\l{}
\l{\sou{estragono}.- Planto herba, aromata, uzata kom kondimento. -}
\l{\cel{3}L.\cel{1}\sou{artemisia dracuncul}es.\cel{2}DEFIRS.}
\l{}
\l{\sou{estribo}. Quaza triangulo de fero, pinta-baza, qua,suspendate}
\l{\cel{3}an singla latero di la selo per rimeno, subtenas la pedo}
\l{\cel{3}di la kavalkanto. - FS.}
\l{}
\l{\sou{estuario}. (\sou{geogr.)}\cel{2}Larja ekflueyo di fluvio, gulfatra. -}
\l{\cel{3}DEFIS.}
\l{}
\l{\sou{estublo}. Parto di la stipo qua restas stacanta kande onu}
\l{\cel{3}falchas la frumento, la sekalo. - DEFI.}
\l{}
\l{\sou{esturjon}o. (\sou{zool.)}\cel{2}Grosa fisho qua iras de la mari a la granda}
\l{\cel{3}fluvii, e di qua la ovi uzesas por facar kaviaro. - L.}
\l{\cel{3}\sou{acipenser Huso}. - EFIS.}
\l{}
\l{\sou{esvanar}. (\sou{netrans.}) Ne plus konciar sua vivo. - FI.}
\l{}
\l{\sou{etajo}. (\sou{arkitekt.}) En domo, en edifico, konsistanta ye plura}
\l{\cel{3}apartmenti superpozita : omna parto sama-nivela. - DEFR.}
\l{}
\l{\sou{etajero}. Moblo ek montanti qui portas tabuleti horizontala,}
\l{\cel{3}dispozita etaje, sur qui onu pozas libri, art-objekti,}
\l{\cel{3}e c. - DFR.}
\l{}
\l{\sou{etalono}. I. Tipo legale di la mezurili e ponderili. - II.}
\l{\cel{3}(\sou{financo}). Metalo selektita por la fabriko di la moneto-}
\l{\cel{3}peco qua uzesas kom tipa unajo di omna moneto-peci di la}
\l{\cel{3}lando koncernata. - F.}
}
